\scanpage{28}
% @ 1–5, top
On considère $X_i=p_i^{*}(X)$ $(i=1,2)$, on va définir un isomorphisme
\[
\deleted{\unsure{\varphi}}\ \psi:\pi^{*}(X_1)\simeq\pi^{*}(X_2)
\]
compatible avec la projection $\pi:W\longrightarrow T\times_S T$, i.e. dont les images inverses sur $W\times_{\pi}W$ \unsure{soient|sont} les \unsure{\^{m}s|m\textsuperscript{es}|mêmes}.

% @ 6–8, cancelled box
\deleted{Pour ceci, on note que}
\[
\begin{aligned}
&\deleted{\pi^{*}(X_1)\simeq p_{\unsure{12|22}}^{(3)*}(X\times_S T|U)}\\
&\deleted{\pi^{*}(X_2)\simeq p_{23}^{(3)*}(T\times_S X|U)}
\end{aligned}
\]

% @ 9–12, middle
Pour ceci, on considère les $\deleted{\unsure{p}}q_i^{(3)*}(X)$ $(i=1,2,3)$,
% slanted
\margin{les $q_i$ induits sur $W$ par les $p_i^{(3)}$}
on note que $\deleted{\unsure{p}}q_1^{(3)*}(X)=\pi^{*}(X_1)$, $\deleted{\unsure{p}}q_3^{(3)*}(X)\simeq\pi^{*}(X_2)$, et qu'on peut considérer le composé
\[
\deleted{\unsure{p}}q_1^{\deleted{\unsure{(3)}}*}(X)\xrightarrow{\sim}\deleted{\unsure{p}}q_2^{\deleted{\unsure{(3)}}*}(X)\xrightarrow{\sim}q_3^{*}(X)
\]

% @ 13–16, lower middle and left margin
Pour prouver la compatibilité de $\psi$ avec $\pi$,
% slanted
\margin{\uline{Question} \uline{de rédaction}}
\deleted{on \unsure{voit|vérifie} $W$}
on écrit $W\times_{T\times_S T}W\subset$
\diagram{\begin{tikzcd}[ampersand replacement=\&,column sep=tiny,row sep=tiny]
T\arrow[d]\&{}\times_S T\arrow[d]\&{}\times_S T\arrow[d]\&{}\times_S T\arrow[d]\\
t_1,\&t',\&t'',\&t_2
\end{tikzcd}}
, et on considère le \deleted{produit} \inserted{sous-objet \deleted{\unsure{ouvert}} $W'$} de
\diagram{\(\begin{array}{ccccccccc}T&\times_S&T&\times_S&T&\times_S&T&\times_S&T\\t_1&&t'&&t''&&t_2&&t'''\end{array}\)}
\deleted{\unsure{formé des septi}} \deleted{\unsure{formé}} correspondant aux \unsure{systèmes|septèmes} $(t_1,t',t'',t_2,t''')$ tels que \deleted{$(t''',t_1),(t''',t_2),(t''',t')$ soient dans $U(\xi)$}
\inserted{$(t''',\deleted{\unsure{t_2}}t',t_1^{\deleted{\unsure{4}}}),(t''',t',t_2)$, \deleted{$(t''',t'')$} $(t''',t'',t_1),(t''',t'',t_2),(t''',t_1,t_2)\in V$.}

% @ lower left, diagram
\diagram{\begin{tikzcd}[ampersand replacement=\&,column sep=small,row sep=small]
\&t'\arrow[dr,"U"]\&\\
t_1\arrow[ur,"U"]\arrow[r,"U"]\arrow[dr,"U"']\&t'''\arrow[u,"U"']\arrow[d,"U"]\arrow[r,"U"]\&t_2\\
\&t''\arrow[ur,"U"']\&
\end{tikzcd}}

% @ final two lines and cancelled formula
On aura donc, si $t_1,t',t'',t_2,t'''$ sont comme il a été spécifié,
\[
\deleted{\varphi_{t_2t'}\varphi_{t't_{2}}=(\varphi_{t_2t'''}\varphi_{t'''t'})(\varphi}
\]
% archive sheet number: 27
